\documentclass[12pt]{article}

\usepackage[utf8]{inputenc}
\usepackage[T1]{fontenc}

\title{LA LINGÜÍSTICA DE LA BIOLOGÍA}
\author{E. CUELLAR, K. HERRERA, K. SANDOVAL, P. SANTIAGO}
\date{\today}

\begin{document}

\maketitle

\section{INTRODUCCIÓN}
Como muchos ya creen, el lenguaje no fue una invención biológica al azar, sino una adptación que surggió mediante la selección natural a lo largo de miles de años. Resulta necesario, por tanto, analizar el momento en que se reunieron las condiciones anatómicas, morfológicas y congnitivas de los \textbf{hominidos}. No debemos olvidar que no existió una única especie de homínidos, sino varias que convergieron evolutivamente y permitieron el surgimiento de la comunicación articulada.
El cambio más importante fue la adopción del bipedismo como forma de dezplazamiento. Caminar sobre dos extremidades no solo liberó nuestras manos para el empleamiento de herramientas, sino que también generó una reestructuracón de las vías respiratorias y del tracto vocal. El descenso de la laringe en la gargata creó una cámara de resonancia más amplia, lo que permitió al \textbf{\textit{Homo sapiens}} modular el flujo de aire y producir una enorme variedad de vocales y consonantes.
Sim embargo, el desarrollo de esta "cámara de resonancia" no fue el único factor evolutivo decisivo. Otros cambios resultaron igual de importantes, especialmente la encefalización progresiva, es decir, el desarrollo del cerebro y el espacio que este ocupa. En nuestra especie, este proceso no se limitó a un simple aumento del tamaño del craneo, sino que impulsó la especialización de áreas cerebrales altamente específicas. 

De hecho, el descenso laríngeo no basta por sí solo para explicar el lenguaje: porque otros mamíferos también los presentan y lo utilizan con fines distintos.Estos mismos factores demuestran que la aparición del lenguaje no puede atribuirse únicamente a modificaciones anatómicas destinadas a la producción de sonidos. La comunicación exige una compleja organización cognitiva. El aumento de la encefalización en los homínidos fue acompañado por mejoras en las capacidades de procesamiento, análisis, memoria y aprendizaje. Estas características sentaron las bases para formas de comunicación cada vez más elaboradas. En este sentido, el lenguaje fue más un producto de la evolución del cerebro que del tracto vocal. 

Aquí entra el pensamiento de Noam Chomsky. Para él, el lenguaje humano no se reduce a un conjunto de sonidos y palabras aprendidos por mera repetición e imitación.(Hauser et al.,2002) El ser humano posee una capacidad innata para contruir y comprender una cantidad casi infinita de expresiones lingüísticas. Desde esta perspectiva, no es necesario aprender cada oración a partir de estímulos externos; basta con adquirir las reglas básicas que permiten generar nuevas estructuras. 

Surge, no obstante, una pregunta central desde la biología evolutiva: ¿como surgió esta capacidad cognitiva y que procesos influyeron en su aparición? El desarollo de habilidades lingüísticas otorgó una ventaja considerable a los organismos con estilos de vida basados en coperación grupal. un individuo capaz de informar a otros miembros del grupo sobre la ubicación del alimento, los peligros del entorno o las estrategias de supervivencia, incrementa su eficacia biológica (fitness) tanto a escala individual como poblacional.

Aun asi, la presión selectiva exacta que favoreció el lenguaje en los hominidos sigue siendo una incógnita. Algunos antropólogos apuntan a las estrategias de caza organizada; otros, a los conflictos con otras especies de homínidos o a la formacion de grupos mas estructurados. todo permanece en suposiciones, pues las primeras formas del lenguaje no dejaron rastos en el registro fósil.

Mientras todas las especies compartimos un código universal, el ADN, los seres humanos poseemos un sistema  de comunicación distinto y organizado de manera jerárquica.Comprender el origen de esta capacidad requiere un enfoque  biolingüistico que integre la lingüistica con la biología evolutiva , la neurociencia y la antropología .Para hablar del tema el problema resulta indispensable distinguir entre la facultad del lenguaje en sentido amplio(FLB) y la facultad del lenguaje del estricto.la FLB engloba la totalidad de los sistemas internos del organismo que participan en el lenguaje , incluidas las interfaces sensorial , motora , conceptual e intencional . La FLB se reduce  al núcleo  de cómputo sintáctico abstracto y especializado que combina elementos. Esta propiedad permite tomar un conjunto finito de unidades discretas y generar variedad ilimitada de estructuras: la llamada infinidad discreta.

Se han formulado tres hipótesis principales al respecto. La primera sostiene que el lenguaje humanon es esencialmente igual al de los animales: no inventamos nada radicalmente nuevo, sino que utilizamos los mismos mecanismos biológicos  y funcionales que otras especies emplean para comunicarse. La segunda afirma que el lenguaje es una adaptación compleja y unica del ser humano , una capacidad sumamente elaborada dentro de los vertebrados. La tercera, finalmente, postula que compartimos con otros animales la FLB; mecanismos de audición, articulación y formación de conceptos heredados de nuestros ancestros pero lo que nos hace unicos e la FLN: solo la recursión nos permite combinar un número finitos de elementos para generar infinitas oraciones.

Resulta necesario examinar el lenguaje no solo como una característica distintiva de nuestra especie, sino como el resultado de una serie de procesos evolutivos, cognitivos y biológicos cuya interacción aún plantea interrogantes. Para comprender el origen es de vitalidad comprender aquellas transformaciones que permitieron la aparicion de sistemas de comunicación cada vez más complejos. De esta manera, se aborda el lenguaje bajo una perspectiva evolutiva en su biología.

Detrás de la reconfiguración del tracto vocal y del aumento de la capacidad craneal existe una historia genética cada vez mejor documentada. Uno de los hallazgos más relevantes en este campo corresponde al gen \textbf{FOXP2}. Este gen codifica un factor de transcripción que regula la expresión de otros genes involucrados en el desarrollo de circuitos neurales relacionados con el control motor fino de la boca, la lengua y la laringe. La versión humana de \textbf{FOXP2} presenta dos sustituciones de aminoacidos que no se encuentran en los parientes primates, tales como el chimpancé. (Ernard et all,. 2002)

Esta evidencia permite considerar que la capacidad lingúistica tiene una base genética; si, aunque no depende de un único componente heredado. por consecuente, resulta más adecuado interpretar la evolución del lenguaje como un proceso de adaptaciones reguladas a las redes del desarrollo neuronal, en lugar de buscar la mutación milagrosa que explique por sí sola el genecis de esta capacidad que consideramos única.

Antes de continuar con el desarrollo evolutivo, resulta necesario definir qué entedemos por lenguaje. En el ámbito social y cultural; el lenguaje constituye una de las principales herramientas mediante las cuales los seres humanos representan su entorno, transmiten conocimientos, expresan pensamientos y construyen relaciones interpersonales. Su importancia no se limita a la producción de sonidos o palabras, puesto que permite organizar las experiencias establecidas a conceptos abstractos y transmitir esas ideas que referencian a las situaciones no presentes que ya se vivieron. De esta manera el lenguaje ha adquirido un papel fundamental en la organización de la sociedad humana por miles de generaciones.

Sin embargo, definir el lenguaje desde la biología requiere mucho más que indentigicarlo como un sistema de comunicación, debido a que diferentes organismos poseen mecanismos especializados para transmitir información mediante señales químicas, acusticas en incluso electricas. La existencia de estas formas de comunicación demuestra que el intercambio de información no es exclusivo de nuestra especie. La diferencia está en la facultad lingúistica humana como un sistema de comunicacion totalemnte diferente a los presentes en el reino animal identificando las carcateristicas unicas del sistema que nos diferencia (Hauser el al., 2002; Cartmill, 2023).

El lenguaje como una capacidad imnata de la biología de los seres vivos, que en caso de la propia raza humana, cuyo desarrollo depende de la interacion entre la estructura neuronal y los sistemas sensoriales que lo estimulan para arrojar una repuesta motriz (a veces). Esta disciplina busca comprender cómo se desarrolla dicha capacidad en los individuos. Por tanto, el lenguaje deja de considerarse únicamente una construcción social y pasa a estudiarse también como un objeto de investigación biologíca (Martins & Boeckx, 2016).

Es relevante como la capacidad combinatoria dle lenguaje humana es peculiar. A partir de un conjunto limitado de unidades, vocablos, sonetos, tonos, silabas, entre otros y las reglas grmaticales que las rigen, los individuos construyen una cantidad practicamente ilimitada de expresiones. Esta propiedad permite organizar elementos en estructuras jerárquicas, producir oraciones cuyo significado es subjetivo sujeto a la interpretacion de quien lo recibe como el que lo emite. Hauser y colaboradores (2002) relacionan esta capacidad con la facultad del lenguaje y plantean una distinción entre sus componentes en sentido amplio (\textbf{FLB}) y en sentido estricto (\textbf{FLN}). La primera comprende los sistemas sensoriales, motores y conceptuales que intervienen en el lenguaje, mientras que la segunda correponde al componente computacional que los autores proponen como posible caracteristica distintiva de nuestra especie.

No obstante, considerar que estas propiedades son exclusivamente humanas requiere examinar con mayor detalle los sistemas de comunicación de otros organismos. Diferentes investigaciones han identificado en animales capacidades relacionadas con el aprendizaje de señales, la comunicación referencial y la organización de secuencias acústicas. Esto no significa que sus sistemas sean equivalentes al lenguaje humano, pero demuestra que algunas de sus características tienen antecedentes o manifestaciones comparables en otras especies. Cartmill (2023) advierte que la comparación entre ambos fenómenos puede conducir a conclusiones equivocas cuando se estudian los animales mediante criterios distintos a los empleados para analizar el habla humana. Por consiguiente, la singularidad del lenguaje humano no debería establecerse únicamente por la presencia o ausencia de caracteristicas aisladas, sino mediante el análisis de la organización y combinación de las diversas capacidades comunicativas.

En este sentido, Bolhuis (2014) y sus colaboradores plantean que una de las propiedades fundamentales que debe explicarse desde la evolución es la capacidad de construir estructuras sintácticas jerárquicas. esta propuesta permite abordar el lenguaje como algo más complejo que un conjunto de señales asociados con determinados estímulos del ambiente. Su organización permite combinar representaciones y establecer relaciones entre ellas, generando expresiones con nuevos significados. Así el problema biológico no consiste solamente en explicar cómo los organismos producen sonidos o reconocen señales, sino en comprender cómo surgió la capacidad cognitiva que permite organizar elementos discretos dentro de estructuras de mayor complejidad.

\section{La comunicación animal como aprocimación al origen biológico del lenguaje}

La comparación entre el lenguaje humano y los sistemas de comunicación de otros animales permite reconocer que algunas de la capacidades relacionadas con la producción y comprensión de señales no aparecieron excusivamente en nuestra especie. En diferentes grupos animales se han identificado mecanismos de comunicacióm que involucran vocalizaciones, movimientos corporales, señales quimicas y estimulos visuales. Estas formas de comunicación complen funciones especificas dentro de sus respectivos contextos ecológicos y sociales, como advertir sobre la presencia de depredadores, localizar alimento, atraer parejas o mantener la cohesión de los grupos. Por esta razó, estudiar la comunicacion animal no significa asumir que todas las especies poseen lenguaje igual que los humanos, sino que identificar cuáles capacidades comunicativas comparten una posible relación evolutiva con aquellas que participan en el lenguaje humano.(Berthet et al., 2023; Suzuki, 2024).

una de las principales diferencias entre la comunicación animal y el lenguaje humano se encuentra en la organización de los elementos que componen los mensajes. En numerosos animales, las señales estan asociadas con situaciones particulares y producen respuestas relativamente especificas en otros individuos.Sin embrago, esto no significa que sus sistemas sean simples o que no tengan buena organizacion. Algunas especies pueden combinar señales, modificaar su comportamiento comunicativo de acuerdo con el contexto y utilizar diferentes modalidades para transmitir informacion. La dificultad consiste en establecer hasta que punto estas combinaciones representan una organización semejante a la estructura lingüística humana. Para ellos, es necesario diferenciar entre la presencia de señales combinadas y la existencia de reglas que permitan generar expresiones cuyo significado dependa de la relación entre sus componentes (Scott-Phillips \& Heintz, 2023).

En los primates no humanos se han documentado comportamientos comunicativos que permiten estudiar algunas de las capacidades que pudieron intervenir en la evolucion del lenguaje. Los grandes simios, por ejemplo, utilizan gestos visuales, tactiles y corporales para dirigir la atencion de otros individuos, solicitar acciones o modificar sus interacciones sociales. La produccion de estos gestos puede depender de la orientacion y atencion del receptor, asi como de la respuesta que este manifieste. Cuando el individuo que emite una sela modifica su comportamiento si no obtiene la respuesta espereda, se evidencia una flexibilidad comunicativa que supera la simplereaccion automatica ante un estimulo. Estas observaciones han llevado a considerar que la simple reaccion automatica ante un estimulo. Estas bservaciones ha llevado a cosiderar que la intecionalidad gestual y la interaccion socila pudieron constituir componentes importantes en los procesos evolutivos que precedieron al lenguaje humano (Prieur et al., 2020).

La importancia de los gestos tambien se relaciona con el caracter multimodal de la comunicacion. Los seres humanos no transmiten informacion exclusivamente mediante sonidos, sino que integran movimientos de las manos, expresiones faciales, direccion de la mirada y postura corporal. De manera semenjante, los primates no humanos combinan vocalizaciones y gestos durante sus interacciones. Esta correspondencia permite considerar que la evolucion del lenguaje pudo involucrar diferentes modalidades de comunicacion, en lugar de depender unicamente de modificaciones anatomicas relacionadas con la produccion vocal. En consecuencia, el estudio de los gestos proporciona informacion sobre las capacidades cognitivas y sociales que pudieron participar en la contruccion de sistemas comunicativos progrevamente mas complejos (Prieur et al., 2020). 

Otro ejemplo relevante se encuentra en las aves que aprenden sus vocalizaciones mediante la interacción con otros individuos.  A diferencia delas especies cuya produccion vocal depende principalmente de patrones innatos, algunas aves canoras adquieren y modifican sus cantos a partir de la exposicion a modelos acusticos. Durante su desarrollo, los individuos escuchan las vocalizaciones de otros miembros de su especie y posteriormente reproducen patrones aprendidos. Este proceso permite determinads caracteristicas del canto se transmitan culturalmente entre generaciones, aunque su variacion permance condicionada por las capacidades biologicas propias de cada especie. La comparacion con el lenguaje humano resulta imortante porque ambos sistemas presentan aprendizaje vocal y transmision cultural, a pesar de que sus mecanismos y resultados comunicativos no son equivalentes (Bruno et al., 2021).

El aprendizaje vocal de las aves demuestra que la capacidad de modificar las vocalizaciones a partir de la experiencia auditiva no constituye una caracteristica exclusiva de los seres humanso. No obstante, aprender sonidos y reproducirlos no implica necesariamente comprenderlos o utilizarlos para construir expresiones con una estructura semantica y sintactica comparable a la de una lengua humana. el canto puede presentar organizacion temporal, patrones recurrentes y variaciones transmitidas socialmente, pero estas propiedades deben analizarse de acuerdo con su funcion y con los mecanismos cognitivos que las producen. por lo tanto, la semejanza entre el aprendizaje vocal de las aves y el desarrollo del habla humana representa una oportunidad para investigar la evolucion de derterminados componenetes biologicos del lenguaje, sin aumir que ambos sistemas poseen la misma capacidad expresiva (Bruno et al., 2021; Scott-Phillips \& Heintz, 2023).

Las investigaciones sobre comunicacion animal tambien han identificado combinaciones de vocalizaciones que parecen transmitir informacion diferente de la que producen las señales individuales. En algunas especies de primates y aves se han observado secuencias de llamadas cuya organizacion modifica la respuesta de los receptores. Estos hallazgos han generado interes porque permiten examinar si existen formas de composicion del significado de una expresion dependa tanto del significado de sus componentes como de la manera en que estos se organizan. En las lenguas humanas, esta propiedad permite construir expresiones nevas mediante la combinacion de unidades lisgüísticas. En los animales, aunque exiten resultados que apuntan hacia combinaciones con significado todavia se discute si los mecanismos identificados corresponden a procesos composicionales equivalentes o a asociaciones especificas entre señales (Suzuki, 2024; Klinedinst, 2025).

Esta discusion permite reconocer la importancia de utilizar criterios rigurosos al comparar los sistemas de comunicacion. No basta con identificar una secuencia de sonidos o movientos para afirmar que una especie posee sintaxis, semantica o lenguaje. Es mecesario establecer que informacion transmite cada señal, como responde el receptor, si el orden de los elementos modofoca el significado y si los individuos pueden producir o comprender combianciones que no hayan experimentado previamente. La lingüística animal ha propuesto precisamente la integracion de comceptos procedentes de la biologia del comportamiento, estas se usan para estudiar las capacidades mediante criterios comparables. Esta aproximacion evita atribuir acracteristicas humanas a los animales unicamente por la semejanza superficial de sus comportamientos y, al mismo tiempo, permite reconocer capacidades comunicativas que tradicionalmente habian sido subestimadas (Berthet et al., 2023).

La comparacion tambien requiere distinguir entre analogia y homologia evolutiva. dos especies pueden presentar comportamientos comunicativos semejante porque ha desarrollado soluciones funcionales similares ante el problemas relacionados con su entorno, aunque estas capacidades no hayan sido heredadas de un antepasado comun que las poseyera. En cambio, una caracteristica homologa corresponde a una capacidad compartida por descendencia evolutiva. Esta distincion es importante porque la presencia de vocalizaciones organizadas de aves y seres humanos no demuestra, por si sola, que ambos sistemas hayan evolucionado a partir de una misma estructura comunicativa ancestral. Las comparaciones entre especies cercanas y distantes permiten estudiar diferentes aspectos del problema: los primates ofrecen informacion sobre capacidades heredadas de antepasados compartidos, mientras que las aves y otros animales permiten examinar la evolucion independiente de mecanismos de aprendizaje, produccion y organizacion de señales (Berthet et al., 2023; Scott-Phillips \& Heintz, 2023).

Por consiguiente, las diferecias entre el lenguaje humano y comunicacion animal no debes establecer a partir de uan oposicion absoluta entre una especie capaz de comunicarse y otras que carecen de esa capacidad. La evidencia disponible muestra que diferentes componentes relacionados con el aprendizaje, la referencia, la intecionalidad y la combiancion de señales aparecen, en distintos grados, en otros organismos. Sin embargo, la presencia de estos componentes por separado no equivale necesariamente a la organizacion integral que caracteriza a las lenguas humanas. La particularidad del lenguaje humano debe examinarse considerando la interaccion entre sus propiedades cognitivas, la organizacion de sus estructuras, su capacidad productiva y las condiciones sociales y culturales en las que se desarrolla. Esta perspectiva permite comprender el lenguaje como un fenomeno biologico que adquirio una dimension cultural extraordinaria, sin separar su evolucion e las capacidades comunicativas presentes en otros seres vivos. 

\section{La evolución del lenguaje: interacción entre biología, cognición y vida social}

Comprender el origen del lenguaje exige considerar que una capacidad biologica no evoluciona de manera independiente de las condiciones ambientales y sociales en las que se desarrolla un organismo. En el caso humano, la comunicacion se encuentra relacionada con la organizacion de los grupos, la cooperacion, el aprendizaje y la transmision de informacion entre individuos. Estas condiciones pudieron favorecer progresivamente sistemas comunicativos mas fexibles, capaces de responder a situaciones que requerian coordinacion y representacion de aconteciminetos que no estaban presentes de manera inmediata. Sin embargo, establecer cual de estos factores tuvo mayor importancia continua siendo un problema abierto, debido a que el origen del lenguaje no puede reconstruirse mediante la observacion directa de las primera poblaciones humanas.

Unas de las explicaciones consideradas en el estudio de la evolucion del lenguaje relaciona el desarrollo de las capacidades comunicativas con el aumento de la complejidad de las interacciones sociales. La vida en un grupos implica reconocer individuos, establecer relaciones, coordinar actividades y responder a comportamientos de otros miembros de la poblacion. En estas circunstancias, la comunicacion puede representar una ventaja al facilitar la cooperacion y disminuir los costos asociados con la coordinacion de acciones. la posibilidad de comparitr informacion sobre recursos, amenazas o actividades colectivas habria contribuido a que los individuos con mejores capacidades de comunicacion obtuvieran ventajas dentro de determinados contextos sociales. No obstante, esta explicacion no debe entenderse como una relacion causal demostrada, sino como una hipotesis que relaciona las presiones sociales con la evolucion de las capacidades cognitivas y comunicativas (Fitch, 2018). 

La cooperacion la cooperacion contituye otro aspecto relevante porque permite examinar por que la comunicacion humana abquirio una capacidad de transmision de informacion tan amplia. A diferencia de los sistemas en los que las señales estan estrechamente relacionadas con estimulos inmediatos, el lenguaje permite comunicar aconteciminetos pasados, anticipar situaciones futuras y compartir informacion sobre elementos que no se encuientran presentes. Esta capacidad amplia las posiblidades de coordinacion entre individuos y permite conservar conocimientos que pueden ser utilizados por otros miembros de la comunidad. De esta manera, el lenguaje no solamente participa en las interacciones inmediatas, sino que tambien interviene en la acimulacion y transmision de conocimentos a travez del tiempo, favoreciendo la continuidad de practicas culturales entre generaciones (Pagel, 2017).

La relacion entre evolucion biologica y transmision cultural permite comprender que el lenguaje no depende exclusivamente de las caracteristicas anatomicas o geneticas de los individuos. Las capacidades biologicas hacen posible abquirir y utilizar una lengua, pero las lenguas particulares se desarrollan y transforman mediante la interaccion entre personas. un individuo nace con capacidades para adquirir lenguaje, pero no nace hablando una lengua especifica. La adquisicion depende de la exposicion a un sistema lingüístico presente en su entorno y de las interacciones que establece durante su desarrollo. Por ellos, la diversidad de lenguas humanas no representan una diversidad de facultades biologicas completamente distintas, sino diferentes expresiones culturales de una capacidad que comparte nuestra especie (Martins \& Boeckx, 2016).

Esta relación entre biologia y cultura tambien permite comprender que el lenguaje modifica las condiciones en las que se desarrolla las propia sociedad humana. la transmisión de conocimientos mediante expresiones lingüísticas posibilita que las experencias individuales se incorporen a un repertorio colectivo, el cual puede ser utilizado y transformado por generaciones posteriores. La enseñanza, la organizacion de actividades, la conservacion de relatos y la construcción de explicaciones sobre el entorno dependen, en gran medida, de esta capacidad. De esta amnera,el lenguaje no debe considerarse únicamente como un resultado de la evolucion humana, sino también como elemento que ha influido en la transformación de las prácticas sociales y culturales de nuestra especie.

\section{Las bases neurobiológicas del lenguaje y su desarrollo}

La capacidad humana para adquirir y utilizar el lenguaje también depende de la organización y el funcionamiento del sistema nervioso. La producción y comprensión de expresiones lingüísticas requieren la participación coordinada de diferentes regiones cerebrales que intervienen en el procesamiento auditivo, la interpretación de significados, la memoria, la planificación motora la integración de información. Por esta razón, no esulta adecuado considerar que el lenguaje se encuentra localizado en una única estructura cerebral. Aunque determinadas regiones presentan una participación importane en algunas operaciones lingüísticas, su funcionamiento depende de redes interconectadas que permiten procesar los distintos componentes de la comunicación (Hickok \& Poeppel, 2007).

Uno de los modelos utilizados para explicar esta organización es el propuesto por Hickok y Poeppel (2007), quienes describen dos grandes vías de procesamiento del habla. La via ventral participa principalmente en la relacion entre las señales acústicas y la comprensión de sus significados, mientras que la vía dorsal interviene en la integración entre la información auditiva y las representaciones necesarias para la produccíon articulatoria. Ambas vías mantienen relaciones funcionales con otras regiones cerebrales, por lo que la comprensión y producción del habla no corresponden a procesos completamente independientes. Esta organización permite interpretar una secuencia de sonidos, relacionarla con conocimientos previos y preparar una respuesta lingüística de acuerto con el contexto. 

La participación de estas redes puede observarse en actividades cotidianas como escuchar una pregunta y formular una respuesta. Inicialmente, el sistema auditivo procesa las características de la señal recibida; posteriormente, diferentes mecanismos permiten reconocer unidades lingüísticas, acceder a sus significados y relacionarlas con la información almacenada en la memoria. Cuando el individuo responde verbalmente, intervienen procesos adicionales de selección de palabras, organización de la expresión y planificación de los movimientos necesarios para producir los sonidos. Aunque estas operaciones pueden parecer inmediatas durante una conversación, requieren una coordinación neuronal considerable. La rapidez con la que se desarrollan demuestra que el lenguaje depende de sismtemas especializados que funcionan de manera integrada y no de una sucesión simple de respuestas aisladas (Hickok \& Poeppel, 2007).

Sin embargo, la organización cerebral relacionada con el lenguaje no aparece completamente desarrollada desde el nacimiento. durante los primeros años de vida, el sistema nervioso experimenta cambios asociados con la maduración biológica y con la interacción del individuo con su entorno. Desde estapas tempranas, los bebés pueden distinguir diferentes características de los sonidos del habla, incluso antes de comprender palabras o producir expresiones reconocibles. Posteriormente, la exposición a una lengua determinada modifica progresivamente la sensibilidad del sistema perceptivo, favoreciendo el reconocimiento de las unidades sonoras que tienen importancia dentro de su entorno lingüístico. Este proceso demuestra que la adquisición del lenguaje requiere tanto una disposición biológica para aprender como experiencias que orienten el desarrollo de las capacidades perceptivas y cognitivas (Kuhl, 2010).

La adquisición de una lengua durante la infancia constituye, por tanto, un proceso en el que interactúan la maduracón cerebral, la experiencia auditiva y la participación social. La exposición a otras personas proporciona información sobre los sonidos, las palabras y las formas de organización que caracterizan una lengua. No obsteante, el aprendizaje no consiste únicamente en imitar las expresiones escuchadas, ya que los niños desarrollan progresivamente la capacidad de reconocer regularidades, establecer relaciones entre palabras y producir expresiones que no han escuchado previamente. Esta capacidad evidencia que el desarrollo lingüístico involucra mecanismos de aprendizaje capaces de extraer patrones a partir de la experiencia y utilizarlos para construir nuevas expresiones (Gervain \& Mehler, 2010).

La relación entre experiencia y desarrollo cerebral también permite abordar la importancia de los períodos sensibles durante la adquisición del lenguaje. A lo largo del desarrollo existen etapas en las que determinados aprendizajes pueden producirse con mayor facilidad debido a la plasticidad de los sistemas nerviosos involucrados. En el procesamiento fonético, por ejemplo, los bebés muestran inicialmente sensibilidad a una amplia variedad de contrastes sonoros presentes en diferentes lenguas. Con el tiempo, la experiencia con la lengua del enotrno contribuye a reorganizar esa sensibilidad, de manera que los contrastes relevantes para su comunicación cotidiana adquieren mayor importancia. Este fenomeno no significa que la capacidad de aprender otras lenguas desaparezca por completo después de la infancia, sino que las condiciones y la facilidad del aprendizaje pueden modificarse conforme avanza el desarrollo (Kuhl, 2004; Werker \& Tees, 2005).

La existencia de estos procesos permite comprender que la adquisición del lenguaje no depende exclusivamente de una programación genética rígida. Los genes participan en la formación y regulación de los sistemas nerviosos que hacen posible el desarrollo ligüístico, pero su actividad se relaciona con procesos de maduración y con las condiciones ambientales.  En consecuencia, el desarrollo de esta capacidad debe analizarse desde una perspectiva que integre diferentes nieveles biológicos, desde la regulación molecular y la formación de circuitos neuronales hasta el comportamiento observable durante las interacciones sociales. Estas relación resulta especialmente importante para la biolingüística, porque permite estudiar cómo una capacidad propia de nuestra especie se desarrolla a partir de mecanismos biológicos que interactúan continuamente con la experiencia (Graham \& Fisher, 2015).

En este contexto, la investigación genética ha permitido identificar asociaciones entre determinados genes y procesos relacionados con el desarrollo del lenguaje. El gen FOXP2 constituye uno de los casos más estudiados debiod a que ciertas alteraciones en su secuencia pueden producir dificultades importantes en el habla y en determinados aspectos del procesamiento lingüístico. Su producto participa en la regulación de la expresión de otros genes y se encuentra relacionado con procesos de desarrollo y funcionamiento de circuitos neuronales. Sin embargo, su participación no permite considerarlo responsable único del lenguaje, puesto que las capacidades lingüísticas dependen de múltiples procesos moleculares, celulares y cognitivos. La investigación sobre FOXP2 demuestra, mas bien, que las alteraciones en determinados componentes de las redes biológicas pueden afectar funciones específicas necesarias para el desarrollo del habla (Fishes, 2019; Graham \& Gisher, 2015).

La interpretación evolutiva de estos hallazgos requiere especial cuidado. La identificación de diferencias genéneticas entre los seres humanos y otros primates no demuestra automáticamente que dichas diferencias hayan originado el lenguaje. Para establecer una relación evolutiva es necesario determinar qué efectos producen las variantes genéticas, en qué tejidos y etapas del desarrollo actúan, cómo modifican los procesos celuares y de qué manera estas modificaciones se relacionan con capacidades cognitivas observables. En el caso de FOXP2, por ejemplo, las investigaciones han revisado la historia evolutiva del gen y las posibles presiones de selección asociadas con determinadas variantes, pero sus resultados no justifican reducir el orgien del leguaje a una modificación genética particular. La evolución de esta capacidad debe entenderse como un proceso complejo en el que pudieron intervenir múltiples cambios genéticos y del desarrollo, junto con transformaciones anatómicas, cognitivas y sociales (Fisher, 2019).

De esta manera, la evidencia neurobiológica y genética permite ampliar la comprensión del lenguaje más allá de su manifestación observable en la conversación. Cada expresión ligüística representa el resultado de procesos que involucran percepción, memoria, aprendizaje, organización conceptual y control motor. Estas operaciones dependen de sistemas biológicos que se desarrollan durante la vida del individuo y que, a su vez, presentan una historia evolutiva propia de nuestra especie. Por consiguiente, estudiar el lenguaje desde la biología implica relacionar las características del comportamiento ligüístico con los mecanismos que lo hacen posible, sin asumir que una sola región cerebral, un gen o una estructura anatómica explica por sí misma la totalidad de esta capacidad.

\section{El lenguaje como capacidad biológica y diversidad cultural}

La relación entre los fundamentos biológicos del lenguaje y la diversidad de las lenguas humanas constituye otro aspecto central para comprender el alcance de la biolinüística. En todas las sociedades humanas conocidas existen existen sistemas ligüísticos que permiten organizar y transmiir la información, aunque sus estructuras gramaticales, sonidos, vocabularios y formas de expresión sean diferentes. Esta diversidad demuestra que el lenguaje no corresponde a un conjunto fijo de palabras o reglas compartidas de manera idéntica por todos los individuos, sino a una capacidad que puede desarrollarse mediante diferentes sistemas convencionales. Las lenguas cambian a lo largo del tiempo, incorporan nuevas expresiones, modifican sus estructuras y responden a las transformaciones históricas de las comunidades que las utilizan (Pagel, 2017).

La diversidad ligüística también permite diferenciar la capacidad humana para adquirir lenguaje de las lenguas particulares que se aprenden durante el desarrollo. El español, el inglés, el francés o las lenguas indígenas de diferentes regiones del mundo no representan capacidades biológicas independientes, sino sistemas culturales que se transmiten mediante la interacción social. Cada uno posee características propias que permiten organizar los sonidos, construir palabras y establecer relaciones gramaticales. A pesar de estas diferencias, los individuos pueden adquirir una lengua cuando cuentan con las condiciones de desarrollo y exposición necesarias. Por ello, el estudio biológico del lenguaje no debe limitarse a describir las características de una lengua específica, sino investigar los mecanismos generales que permiten a los seres humanos aprender y utilizar cualquiera de ella (Martins \& Boeckx, 2016).

Esta capacidad tampoco se encuentra restringida a la comunicación oral. las lenguas de señas constituyen sistemas ligüísticso completos que utilizan configuraciones de las manos, movimientos, expresiones faciales y relaciones espaciales para construir significados. Su existencia permite reconocer que el lenguaje no depende necesariamente de la producción de sonidos mediante el aparato fonador, sino de la capacidad para organizar unidades simbólicas y establecer relaciones entre ellas. Además, las lenguas de señas presentan estructuras gramaticales propias y permiten expresar contenidos abtractos, narrar acontecimientos y desarrollar formas complejas de interacción. Por consiguiente, cualquier explicación biológica del lenguaje debe distinguir entre la capacidad lingüística y una de sus posibles modalidades de expresión, como es el habla (Sandler \& Lillo-Martin, 2006).

El reconocimiento de las lenguas de señas también tiene consecuencias para el estudio de la evolución humana. Si el lenguaje puede manifestarse mediante modalidades visuales y gestuales, su origen no puede explicarse exclusivamente a partir de la evolución de la laringe o de las modificaciones anatómicas relacionadas con la vocalización. Estas estructuras pudieron desempeñar un papel importante en la consolidación del habla, pero no representan por sí solas la totalidad de las capacidades lingüísticas. La existencia de sistemas lingüísticos visuales demuestra que la organización simbólica, la combinación de unidades y la construcción de significados pueden desarrollarse mediante canales sensoriales y motores diferentes. En consecuencia, la evolución del lenguaje debe examinarse considerando tanto los mecanismos cognitivos generales como las adaptaciones que hicieron posible distintas formas de expresión (Sandles \& Lilo-Martin, 2006).

La diversidad de las lenguajes también permite reconocer que el lenguaje participa activamente en la construcción de las identidades y las relaciones sociales. Mediante él se transmiten conocimientos sobre el entorno, prácticas culturales, explicaciones del mundo y formas particules de organizar la experiencia. Una comunidad no utiliza su lengua únicamente para intercambiar información, sino también para conservar parte de su historia y establecer vínculos entre sus integrantes. En este sentido, el lenguaje adquiere una dimensión que supera su función comunicativa inmediata, puesto que participa en la continuidad de los conocimientos y las prácticas que caracterizan a diferentes grupos humanos. Esta dimensión cultural no reemplaza sus fundamentos biológicos, sino que representa uno de los ámbitos en los que dicha capacidad se manifiesta y adquiere un significado (Pagel, 2017).

Por esta razón, la relación entre la biología y la cultura no debe plantearse como una oposicíon en la que una dimensión excluya a la otra. La capacidad de adquirir lenguaje depende de condiciones biológicas, pero su desarrollo ocurre en entornos sociales donde los individuos aprenden sistemas lingüísticos específicos. A su vez, las lenguas transmitidas culturalmente permiten modificar las formas de interacción, conservar conocimientos y generar nuevas prácticas colectivas. Se establece asi una relación continua entre las características del organismo y las condiciones de la vida de las poblaciones humanas. Esta interacción constituye uno de los aspectos que hacen del lenguaje un objeto de estudio particularmente relevante para la biología, porque permite examinar cómo una capacidad desarrollada mediante procesos evolutivos participa posteriormente en la transformación del comportamiento y de la organización social de nuestra especie.

En vistas de todo lo anterior, queda claro que el lenguaje humano constituye una capacidad biológica cuya complejidad no puede explicarse mediante un único factor evolutivo,anatómico,genético o cognitivo. Su desarollo  es el resultado de la interracción de diferentes  procesos que involucran la organización  del sistema nervioso, las capacidades de aprendizaje, la evolución de estructuras relacionadas con la comunicación  y las condiciones  sociales en las que se desenvolvieron nuestros antepasados. Por esta razón, comprender el lenguaje exige superar las explicaciones  que lo reducen  exclusivamente al habla o la producción de sonidos, y reconocerlo como una capacidad que integra procesos biológicos y cognitivos  para construir y transmitir significados.

La comparación con otros animales permite reconocer que el lenguaje humano no surgió en un escenario completamente separado de las capacidades comunicativas presentes en la naturaleza. Diferentes especies presentan mecanismos de aprendizaje vocal, comunicación gestual, transmisión de señales y combinación de elementos comunicativos que permiten estudiar posibles antecedentes o componentes relacionados con nuestra capacidad lingÜistica. Sin embargo estas semenjanzas no significan que todos los sitemas de comunicación sean equivalentes. La particularidad humana se encuentra en la organización e integración de múltiples capacidades que permiten construir expresiones productivas, transmitir  información  compleja y desarollar sistemas simbólicos compartidos por comunidades enteras. En consecuencia estudiar la comunicación animal no dismunuye la singularidad del lenguaje humano, sino que permite comprender mejor su evolución y reconocer la diversidad de soluciones  comunicativas desarolladas por los organismos.
Asimismo, l.os avances en genética y neurobiológía  han demostrado que el lenguaje depende de mecanismo biologícos distribuidos y relacionados entre si. La participación de genes FOXP2, la organización de redes neuronales  y los procesos de plasticidad  durante el desarollo evidencian que esta capacidad  tiene fundamentos materiales suceptibles de investigación científica. No obstante, ninguno de estos elementos  explican por si  solo la totalidad del lenguaje. las diferencias  genéticas, las estructuras cerebrales  y las características anatómicas deben interepretarse en relación  con los procesos de desarollo  y las funciones que desempeñan. De esta manera, la biolingÜistica  permite establecer  conexiones entre mismos niveles de organización biológica sin reducir un fenómeno complejo a una sola estructura o mecanismo.
Por otra parte, la relación entre las bases biológicas y la diversidad cultural demuestra que el lenguaje no puede commprenderse únicamente como una adaptación que permitió  intercambiar información. Su capacidad para conservar conocimientos,  transmitir  experiencias experiencias y organizar  relaciones sociales  ha contribuido a la tranformación  de las comunidades humanas. Las lenguas de señas demuestranque dicha capacidad no depende de la vocalización. Estos aspestos permiten reconocer que la biología establece las condiciones para el desarollo del lenguaje, pero son las experiencias  y las interacciones  sociales las que participan en la adquisición  y transformación de sus diferengtes expresiones.
En definitiva, la biolingÜistica ofrece una perspectiva integrado5ra para estudiar el lenguaje como parte de la historia evolutiva  y del desarollo biológico humano. Su importacia radica en relacionar los mecanismos genéticos, neuronales, cognitivos y sociales que intervienen en esta capacidad, evitando interpretaciones que separen artificialmenteal organismo y su entorno. El lenguaje representa, por lo tanto una de las expresiones más complejas de la interacción entre la evolución biológica y cultura: surgió y se desarrollo  sobre las capacidades propias de los organismos, pero alcanzó en los seres humanos una dimensión simbólica y social que tranformó profundamente sus formas en relacionarse, aprender y construir  conocimiento. Estudiarlo desde la biología  no significa limitarlo  a sus componentes orgánicos, sino comprender cómo estos hacen posible una capacidad que continúa definiendo buena parte de la experiencia humana.
...

\section{BIBLIOGRAFÍA}
\begin{itemize}
    \item Berthet, M., Coye, C., Dezecache, G., \& Kuhn, J. (2023). Animal linguistics: A primer. \textit{Biological Reviews, 98}(1), 81--98. https://doi.org/10.1111/brv.12897
    \item Bruno, J. H., Jarvis, E. D., Liberman, M., \& Tchernichovski, O. (2021). Birdsong learning and culture: Analogies with human spoken language. \textit{Annual Review of Linguistics, 7}, 449--472. https://doi.org/10.1146/annurev-linguistics-090420-121034
    \item Enard, W., Przeworski, M., Fisher, S. E., Lai, C. S. L., Wiebe, V., Kitano, T., Monaco, A. P., \& Pääbo, S. (2002). Molecular evolution of FOXP2, a gene involved in speech and language. \textit{Nature}, \textit{418}(6900), 869–872. https://doi.org/10.1038/nature01025
    \item Fitch, W. T. (2018). The biology and evolution of language: “Deep homology” and the “Sapir-Whorf hypothesis” revisited. \textit{Current Opinion in Behavioral Sciences, 21}, 1--6. https://doi.org/10.1016/j.cobeha.2017.12.003
    \item Fisher, S. E. (2019). Human genetics: The evolving story of FOXP2. \textit{Current Biology, 29}(2), R65--R67. https://doi.org/10.1016/j.cub.2018.11.047
    \item Gervain, J., \& Mehler, J. (2010). Speech perception and language acquisition in the first year of life. \textit{Annual Review of Psychology, 61}, 191--218. https://doi.org/10.1146/annurev.psych.093008.100408
    \item Graham, S. A., \& Fisher, S. E. (2015). Understanding language from a genomic perspective. \textit{Annual Review of Genetics, 49}, 131--160. https://doi.org/10.1146/annurev-genet-120213-092236
    \item Hauser, M. D., Chomsky, N., \& Fitch, W. T. (2002). The faculty of language: What is it, who has it, and how did it evolve? \textit{Science}, \textit{298}(5598), 1569-1579. https://doi.org/10.1126/science.298.5598.1569
    \item Hickok, G., \& Poeppel, D. (2007). The cortical organization of speech processing. \textit{Nature Reviews Neuroscience, 8}, 393--402. https://doi.org/10.1038/nrn2113
    \item Klinedinst, N. (2025). Compositionality in human and animal communication. \textit{Animal Cognition, 28}, Article 90. https://doi.org/10.1007/s10071-025-02010-8
    \item Kuhl, P. K. (2004). Early language acquisition: Cracking the speech code. \textit{Nature Reviews Neuroscience, 5}, 831--843. https://doi.org/10.1038/nrn1533
    \item Kuhl, P. K. (2010). Brain mechanisms in early language acquisition. \textit{Neuron, 67}(5), 713--727. https://doi.org/10.1016/j.neuron.2010.08.038
    \item Martins, P. T., \& Boeckx, C. (2016). What we talk about when we talk about biolinguistics. \textit{Linguistics Vanguard, 2}(1), 20160007. https://doi.org/10.1515/lingvan-2016-0007
    \item Pagel, M. (2017). Q\&A: What is human language, when did it evolve and why should we care? \textit{BMC Biology, 15}, Article 64. https://doi.org/10.1186/s12915-017-0405-3
    \item Prieur, J., Barbu, S., Blois-Heulin, C., \& Lemasson, A. (2020). The origins of gestures and language: History, current advances and proposed theories. \textit{Biological Reviews, 95}(3), 531--554. https://doi.org/10.1111/brv.12576
    \item Sandler, W., \& Lillo-Martin, D. (2006). \textit{Sign language and linguistic universals}. Cambridge University Press. https://doi.org/10.1017/CBO9781139163910
    \item Scott-Phillips, T., \& Heintz, C. (2023). Animal communication in linguistic and cognitive perspective. \textit{Annual Review of Linguistics, 9}, 93--111. https://doi.org/10.1146/annurev-linguistics-030421-061233
    \item Suzuki, T. N. (2024). Animal linguistics. \textit{Annual Review of Ecology, Evolution, and Systematics, 55}, 205--226. https://doi.org/10.1146/annurev-ecolsys-102622-030253
    \item Werker, J. F., \& Tees, R. C. (2005). Speech perception as a window for understanding plasticity and commitment in language systems of the brain. \textit{Developmental Psychobiology, 46}(3), 233--251. https://doi.org/10.1002/dev.20060

\end{itemize}


\end{document}
